\documentclass[12pt]{article}
%\documentclass{article}
\usepackage{ xcolor}
%\textheight 9in
%\textwidth 6in
%\topmargin 0in
%\oddsidemargin 0.5in
%\evensidemargin 0.5in

\title{Laboratorio 1: scrivere formule matematiche}


\begin{document}
%\date{}
\maketitle


Questa \`e una linea di                 
 testo scritta in \LaTeX.










\bigskip

Due matrici $n \times n$ complesse
$A$ e $B$ si dicono \emph{simili}
se esiste una matrice $n \times n$
invertibile $T$ tale che
$$
B=T^{j_{-1}}AT
$$
Questa \`e una formula matematica centrata in un rigo:
\[
    \sum_{i=1}^n \frac{2^3}{1+i}
\]


Questa formula invece resta dentro il rigo del testo
$\sin(2x) = 2 \sin(x) \cos(x) $ mentre si illustra la teoria.


\medskip

Dato il polinomio $P(x)=x^2-2x+\frac 34$, allora
$$
\frac {P(x)}{x-\frac12}=x-\frac 32  .
$$
Cosa si pu\`o dire delle radici di $P$?

\medskip
\hrule

\[
\frac{\sqrt{a^2+b^2}}{\frac{b}{i+1}}
\]
\end{document}

